\section{前缀码与 Kraft 不等式}
\label{sec:07-Information-Theory/04-Source-Coding/01}

三个符号 \(A,B,C\)，概率分别是 \(0.5,0.25,0.25\)。要压缩，直觉很明确：常见的给短码字。于是写下

\[
A\mapsto 0,\qquad B\mapsto 01,\qquad C\mapsto 1.
\]

三个码字互不相同，映射是单射，看上去挑不出毛病。现在把比特串 \(01\) 交给译码器。它读到第一个 \(0\)，这已经是 \(A\) 的完整码字，于是输出 \(A\)，再把剩下的 \(1\) 读成 \(C\)。可它同样有理由再等一个比特，把 \(01\) 整个读成 \(B\)。一串比特两种拆法，发送方的原意再也找不回来。

问题出在拼接。编码器把码字首尾相连发出去，中间不插分隔符，而分隔符恰恰是译码器最想要的东西。所以在比较``哪套码更短''之前，得先划清``哪套码根本能用''。

\subsection{可译性分三层，实际只用最强的那层}

最弱的一层叫非奇异码，只要求不同符号配不同码字。上面那套码满足它，照样崩了。中间一层叫唯一可译码：任何有限长的比特串都只有一种符号分解方式。这才是真正需要的性质，可它不好直接检验——待检的串没有长度上限，穷举无从下手。最强的一层是前缀码，要求任何码字都不是另一个码字的前缀。

前缀码是唯一可译码的真子集，但换来两样好处。检验只需两两比对码字，一眼可见。译码可以边读边输出：沿着码树往下走，走到某个码字的末尾立刻吐出符号并跳回树根，不必前看，不必缓存，不必回溯。

限定在前缀码里，会不会为了这份便利牺牲压缩率？不会。稍后 Kraft--McMillan 不等式会保证：任何唯一可译码的码长组合，都能被某个前缀码原样复制。平均码长只由码长组合决定，所以讨论最优码时把视野收在前缀码内，一分钱都不亏。

\subsection{一个码字吃掉一整棵子树}

把所有二元串摆到一棵无限二叉树上。根是空串，向左走一步添 \(0\)，向右走一步添 \(1\)，深度 \(\ell\) 的节点就是长度 \(\ell\) 的串。选中一个节点作码字，等于宣布它的全部后代作废——那些串统统以这个码字开头，一旦启用就会撞车。

深度因此成了价格。深度 \(\ell\) 的节点，其后代在整棵树中占的份额恰好是 \(2^{-\ell}\)：短码字贵，长码字便宜。前缀性要求各码字的子树互不相交，而所有子树都塞在同一棵树里，份额之和不能超过一。

\begin{center}
\begin{tikzpicture}[x=1cm,y=1cm,font=\small,>=stealth]
  \fill[blue!14] (-1.5,1.7) -- (-2.7,-0.5) -- (-0.3,-0.5) -- cycle;
  \fill[blue!14] (0.7,0.6) -- (0.1,-0.5) -- (1.3,-0.5) -- cycle;
  \fill[blue!14] (2.3,0.6) -- (1.7,-0.5) -- (2.9,-0.5) -- cycle;
  \draw[black!55] (0,2.8) -- (-1.5,1.7);
  \draw[black!55] (0,2.8) -- (1.5,1.7);
  \draw[black!55] (1.5,1.7) -- (0.7,0.6);
  \draw[black!55] (1.5,1.7) -- (2.3,0.6);
  \node[font=\footnotesize,black!70] at (-0.98,2.45) {\(0\)};
  \node[font=\footnotesize,black!70] at (0.98,2.45) {\(1\)};
  \node[font=\footnotesize,black!70] at (0.82,1.32) {\(0\)};
  \node[font=\footnotesize,black!70] at (2.18,1.32) {\(1\)};
  \fill[black!70] (0,2.8) circle (1.7pt);
  \fill[black!70] (1.5,1.7) circle (1.7pt);
  \fill[red!75] (-1.5,1.7) circle (2.3pt);
  \fill[red!75] (0.7,0.6) circle (2.3pt);
  \fill[red!75] (2.3,0.6) circle (2.3pt);
  \node[above,font=\footnotesize,black!60] at (0,2.88) {根（空串）};
  \node[font=\footnotesize] at (-1.5,-0.9) {\(A\mapsto 0\)};
  \node[font=\footnotesize] at (0.7,-0.9) {\(B\mapsto 10\)};
  \node[font=\footnotesize] at (2.3,-0.9) {\(C\mapsto 11\)};
  \node[font=\footnotesize,black!60] at (-1.5,-1.4) {锁掉 \(1/2\)};
  \node[font=\footnotesize,black!60] at (0.7,-1.4) {锁掉 \(1/4\)};
  \node[font=\footnotesize,black!60] at (2.3,-1.4) {锁掉 \(1/4\)};
  \draw[black!25] (3.7,-1.7) -- (3.7,2.7);
  % 预算条
  \fill[blue!28] (4.6,0.9) rectangle (7.1,1.6);
  \fill[blue!16] (7.1,0.9) rectangle (8.35,1.6);
  \fill[blue!16] (8.35,0.9) rectangle (9.6,1.6);
  \draw[black!50] (4.6,0.9) rectangle (9.6,1.6);
  \draw[black!50] (7.1,0.9) -- (7.1,1.6);
  \draw[black!50] (8.35,0.9) -- (8.35,1.6);
  \node[font=\footnotesize] at (5.85,1.25) {\(2^{-1}\)};
  \node[font=\footnotesize] at (7.72,1.25) {\(2^{-2}\)};
  \node[font=\footnotesize] at (8.97,1.25) {\(2^{-2}\)};
  \node[above,font=\footnotesize,align=center] at (7.1,1.72) {码长 \((1,2,2)\)：和为 \(1\)，恰好装满};
  \fill[blue!28] (4.6,-0.7) rectangle (7.1,0.0);
  \fill[blue!28] (7.1,-0.7) rectangle (9.6,0.0);
  \fill[red!25] (9.6,-0.7) rectangle (10.85,0.0);
  \draw[black!50] (4.6,-0.7) rectangle (9.6,0.0);
  \draw[red!70] (9.6,-0.7) rectangle (10.85,0.0);
  \node[font=\footnotesize] at (5.85,-0.35) {\(2^{-1}\)};
  \node[font=\footnotesize] at (8.35,-0.35) {\(2^{-1}\)};
  \node[font=\footnotesize,black!70] at (10.22,-0.35) {\(2^{-2}\)};
  \node[below,font=\footnotesize,align=center] at (7.4,-0.78) {码长 \((1,1,2)\)：和为 \(5/4\)，第三个码字没地方放};
  \draw[red!70,dashed] (9.6,-0.82) -- (9.6,2.3);
  \node[above,font=\footnotesize,red!70] at (9.6,2.32) {预算上限 \(1\)};
\end{tikzpicture}

\emph{左：选一个节点当码字，就把它下面整棵子树锁死，深度 \(\ell\) 的节点锁掉的份额是 \(2^{-\ell}\)。右：把这些份额并排放进一条长度为 \(1\) 的预算条，Kraft 不等式说的正是它们不能重叠、也不能溢出。}
\end{center}

\begin{theorem}
正整数 \(\ell_1,\dots,\ell_m\) 能成为某个二元前缀码的码长，当且仅当
\[
\sum_{i=1}^m 2^{-\ell_i}\le 1.
\]
码字母表有 \(D\) 个符号时，把底数 \(2\) 换成 \(D\)。
\end{theorem}

必要性就是上一段那句话：子树互不相交，份额加起来装得下。充分性是一个构造。把码长从小到大排序，依次在树上挑一个尚未被封锁的节点。轮到第 \(k\) 个码字时，前面的码字总共占掉 \(\sum_{i<k}2^{-\ell_i}\) 的份额，按不等式这个数严格小于 \(1\)，剩下的空隙至少还有 \(2^{-\ell_k}\)，深度 \(\ell_k\) 上必定还有位置。短码先挑，因为它的候选位置最少，晚了就被挤没了。

McMillan 把这个约束推到了唯一可译码：唯一可译这个更弱的条件，也强制同一个不等式。于是长度组合 \((1,1,2)\) 不但没有前缀码，连唯一可译码都不存在——两个长度为 \(1\) 的码字已经把根的两条分支占满，第三个码字无处可去。

有一件事不等式不管：它只看长度，不看字符串怎么摆。长度 \((1,2,2)\) 通过检验，但集合 \(\set{0,01,11}\) 依然有前缀冲突，\(0\) 是 \(01\) 的前缀。预算不超支只是能实现的前提，钱具体怎么花还得自己核对。

\subsection{码长就是概率的另一种写法}

不等式还有一个读法，比``能不能构造''重要得多。若一组码长使 Kraft 和恰为 \(1\)，令 \(q_i=2^{-\ell_i}\)，这组数非负、求和为一，它是一个概率分布。反过来，任给一个分布 \(q\)，取 \(\ell_i=-\log_2 q_i\) 就把预算不多不少地用尽。两个方向都通。

\begin{quote}
一套前缀码和一个概率分布是同一件事的两种写法，靠 \(\ell_i\leftrightarrow-\log_2 q_i\) 换算。设计编码就是分配概率，指定概率就是指定码长。
\end{quote}

这条对应决定了本章后面所有内容的形状。它也解释了一条在深度学习里天天走的流水线：网络输出 logits，softmax 把它们变成概率，交叉熵损失把概率变成对数代价——而它的单位就是比特。训练一个语言模型和训练一个压缩器，在这条链上是同一个动作，区别只在最后有没有真的把比特写进文件。

回到编码。若符号出现概率是 \(p_i\)，最省的码长分配显然是让 \(q\) 对上 \(p\)：

\[
\ell_i^{*}=-\log_2 p_i.
\]

这组理想长度有两个好性质。按概率加权求平均，得到的正是熵 \(\sum_i p_i\ell_i^{*}=H(X)\)；同时 \(2^{-\ell_i^{*}}=p_i\)，Kraft 和恰为 \(1\)，从不违规。它唯一的毛病是几乎从不是整数：\(p_i=0.3\) 时 \(\ell_i^{*}\approx1.74\)，而你没法发出 \(1.74\) 个比特。

最省事的补救是向上取整，\(\ell_i=\lceil-\log_2 p_i\rceil\)，这叫 Shannon 码。取整只会让 \(2^{-\ell_i}\le p_i\)，Kraft 和更不会超，前缀码一定存在；同时每个码长超出理想值不到一个比特。两边按 \(p_i\) 加权求和，得到

\[
H(X)\le L<H(X)+1,
\]

其中 \(L=\sum_i p_i\ell_i\)。右边是构造给的，说明熵这条线总能贴到不到一比特之内。左边还欠一个证明——凭什么任何码都跨不过去。

\subsection{跨不过熵，是因为编码就是在押一个分布}

取任意一组满足 Kraft 不等式的码长，记

\[
K=\sum_i 2^{-\ell_i}\le 1,\qquad q_i=\frac{2^{-\ell_i}}{K}.
\]

\(q\) 非负且求和为一，是个正经的概率分布——它就是这套码在暗中押注的那个分布。反解码长得 \(\ell_i=-\log_2 q_i-\log_2 K\)，代进平均码长：

\[
\begin{aligned}
L&=\sum_i p_i\ell_i\\
 &=-\sum_i p_i\log_2 q_i-\log_2 K\\
 &=H(P)+\KL(P\|Q)-\log_2 K\\
 &\ge H(P).
\end{aligned}
\]

第三个等号是交叉熵的分解，来自\hyperref[sec:07-Information-Theory/03-Divergences/01]{``错误模型的代价：从交叉熵到 KL 散度''}。最后一步同时用掉两个非负性：散度非负，以及 \(K\le1\) 使 \(-\log_2 K\ge0\)。

这个式子把下界的来源说破了。编码从来不是在跟熵较劲，它是在押一个分布 \(Q\)；押对了，平均码长就是熵，押错了，多花的比特分毫不差地等于 \(\KL(P\|Q)\)。至于 \(K<1\) 那一项，是码树留了空位没占——空着也要付租金。

顺带一个推论值得记住：压缩率是模型质量的直接读数。同一份数据交给两个模型压，谁压得短谁的分布更接近真相，差出来的比特数就是两者散度之差。这解释了为什么可以用压缩率给模型排名，也解释了这种排名的边界——它测的是分布拟合得准不准，不测别的。

\subsection{什么时候取等，什么时候还有余地}

等号需要两件事同时发生：\(P=Q\) 且 \(K=1\)。翻译过来，就是每个概率都恰好是 \(2^{-k}\) 的形式。此时理想码长本来就是整数，取整不损失分毫，Huffman 与 Shannon 码都能压到熵上。

概率一般不长这样，于是逐符号编码总要付一笔``找零''：不到一个比特，但摊到每个符号头上。符号越少、分布越偏，这笔零头占比越难看。极端一点，Bernoulli\((0.99)\) 的熵只有 \(0.081\) 比特，而二元单符号码最短也得给每个符号 \(1\) 比特——差了一个数量级，且不是设计不够聪明造成的。

那不到一比特的差额有两条出路，后面几节各走一条：把多个符号打包成块，让整数取整的浪费摊薄；或者干脆放弃逐符号发码，把整段序列映到一个区间上，让分数比特在序列内部自由累积。

Shannon 码只做到了``可行''，还没做到``最短''。给定一组概率，究竟哪一组整数码长使平均长度最小？下一节给出答案，并说明为什么一个只看最小两个概率的局部动作，足以保证全局最优。
